
\textbf{Théorème 5.3.2.} \uline{L'%
application} \(c_1\) (5.3.1) \uline{est un homomorphisme injectif, et l'homomorphisme}
\[
 \det:\Filt^1(X)\longrightarrow\Pic(X)
\]
\uline{induit, par passage au quotient, un homomorphisme inverse à gauche de} \(c_1\) :
\[
 \alpha:\Gr^1(X)\longrightarrow\Pic(X).
\]
\uline{Si toute section} \(\underline{n}\) \uline{du faisceau constant} \(\mathbb Z_X\) \uline{est bornée (par exemple si} \(X\) \uline{est «~quasi-compact~», ou si l'objet final} \(\underline{e}\) \uline{de} \(X\) \uline{est somme finie d'objets connexes), alors} \(c_1\) \uline{est un isomorphisme et} \(\alpha\) \uline{son inverse.}

\emph{Démonstration} : \(c_1\) est un homomorphisme, car on a pour tout \(L,L'\in\Pic(X)\) l'équation
\[
 (1-(LL')^{-1})-(1-L^{-1})-(1-{L'}^{-1})
 =
 -(1-L^{-1})(1-{L'}^{-1})\in\Filt^2(X).
\]

Admettons pour l'instant que
\[
 \det(\Filt^2(X))=\{1\}\in\Pic(X)
\]
i.e. que \(\alpha\) soit bien défini. On a alors
\[
 \alpha\circ c_1=\id_{\Pic(X)},
\]
car par 5.2.2 b) on a :
\[
 \det(1-L^{-1})=(\det(L-1))=L
\]
pour tout \(L\in\Pic(X)\). Puis on montre que
\[
 c_1\circ\alpha=\id_{\Gr^1(X)}
\]
lorsque toute section \(\underline{n}\) de \(\mathbb Z_X\) est bornée. Comme \(\Filt^1(X)\) est le groupe abélien engendré par les éléments \(\cl^{\bullet}(\underline{E})-\underline{n}\), où \(\underline{E}\) est un
